\section*{Multivariable Calculus}

\subsection*{Basic Types of Functions}

\begin{enumerate}

\item What is a scalar-valued function of several variables? Give its type and notation.

\item What is a vector-valued function of one variable? Give its type and component form.

\item What is a vector-valued function of several variables? Give its type and component form.

\item What is a scalar field? What is a vector field?

\item When can two functions be composed? State the condition in terms of domains and codomains.

\item Give examples of composing:
(a) a scalar-valued function with a vector-valued function,
(b) two vector-valued functions.

\item What is the difference between a coordinate transformation and a vector field?

\end{enumerate}

\subsection*{Vector-Valued Functions of One Variable (Curves)}

\begin{enumerate}

\item How do we define the limit of a vector-valued function $\mathbf{r}(t)$?

\item How do we define continuity of a vector-valued function?

\item How do we define the derivative of $\mathbf{r}(t)$? Give the definition from first principles.

\item What is the geometric meaning of $\mathbf{r}'(t)$?

\item What is the equation of the tangent line to a curve at $t_0$?

\item State the differentiation rules (sum, scalar multiplication, dot product, cross product).

\item Prove that if $|\mathbf{r}(t)| = c$, then $\mathbf{r}(t) \cdot \mathbf{r}'(t) = 0$.

\item How do we define the integral of a vector-valued function?

\item Define the arc length of a curve.

\item What does it mean to parameterize a curve by arc length?

\item What is the unit tangent vector $\mathbf{T}$?

\item What is the unit normal vector $\mathbf{N}$? Why is it orthogonal to $\mathbf{T}$?

\item What is the unit binormal vector $\mathbf{B}$?

\item What is the osculating plane?

\item What is the normal plane?

\item What is curvature? Give its definitions.

\item What is the curvature of a line? Of a circle?

\item What are position, velocity, acceleration, and speed?

\item Decompose acceleration into tangential and normal components.

\item Give examples of purely tangential and purely normal acceleration.

\item What is torsion? What does it measure?

\item What are the units of curvature and torsion?

\end{enumerate}

\subsection*{Functions of Several Variables (Scalar Fields)}

\begin{enumerate}

\item What is a level set of $f:\mathbb{R}^n \to \mathbb{R}$?

\item What is a level curve? What is a level surface?

\item How can level surfaces be reduced to level curves?

\item How do we define the limit of $f:\mathbb{R}^n \to \mathbb{R}$?

\item What does it mean that a limit is independent of path?

\item How do we define continuity?

\item What is a partial derivative? Give the definition.

\item What is the geometric interpretation of partial derivatives?

\item What does it mean for a function to be differentiable?

\item What is the total differential?

\item What is the linear approximation of a function?

\item Describe composition with other functions and with a path $\mathbf{r}(t)$.

\item State the chain rule (general and along a path).

\item What are independent and dependent variables in implicit equations?

\item Given $F(x,y)=0$, derive $\frac{dy}{dx}$ and $\frac{dx}{dy}$.

\item What is the relationship between $\frac{dy}{dx}$ and $\frac{dx}{dy}$?

\end{enumerate}

\subsection*{Gradients, Geometry, and Optimization}

\begin{enumerate}

\item What is the gradient $\nabla f$? Give its type and notation.

\item What is the geometric interpretation of the gradient?

\item Why is the gradient orthogonal to level curves? Prove it.

\item Why is the gradient orthogonal to level surfaces? Prove it.

\item What is the directional derivative?

\item How is the directional derivative related to the gradient?

\item In which direction is the directional derivative maximal?

\item What is the derivative of $f$ along a path?

\item What is a critical point?

\item State Fermat's theorem.

\item What is the second derivative test?

\item What is the equation of a tangent line to a level curve?

\item What is the equation of a tangent plane to a level surface?

\item What is the equation of a normal line?

\item What is the equation of the tangent plane to $z=f(x,y)$?

\end{enumerate}

\subsection*{Implicit Functions and Constrained Optimization}

\begin{enumerate}

\item What problem does the Implicit Function Theorem solve?

\item State the Implicit Function Theorem for $F(x,y)=0$.

\item What condition ensures that $y$ can be written as a function of $x$?

\item Derive $\frac{dy}{dx} = -\frac{F_x}{F_y}$.

\item What is the geometric meaning of $F_y \neq 0$?

\item What happens if $F_y = 0$?

\item Extend the theorem to $F(x,y,z)=0$.

\item Express $\frac{\partial z}{\partial x}$ and $\frac{\partial z}{\partial y}$.

\item How does the theorem justify implicit differentiation?

\item Give examples where the theorem applies and fails.

\item What is the problem of constrained optimization?

\item State the method of Lagrange multipliers.

\item What system of equations must be solved?

\item Why must $\nabla f = \lambda \nabla g$?

\item What is the geometric interpretation of this condition?

\item How does the method extend to higher dimensions?

\item What happens if $\nabla g = \mathbf{0}$?

\end{enumerate}

\subsection*{Vector-Valued Functions of Several Variables and Jacobians}

\begin{enumerate}

\item What is a vector-valued function $F:\mathbb{R}^n \to \mathbb{R}^m$?

\item What is a vector field?

\item What is a parametric surface?

\item How do we compute tangent vectors to a parametric surface?

\item What is the tangent plane to a parametric surface?

\item What is the Jacobian matrix?

\item How is it defined in terms of partial derivatives?

\item What do the rows represent?

\item What do the columns represent?

\item When does the Jacobian reduce to the gradient?

\item What is the total derivative?

\item How does the Jacobian represent linear approximation?

\item What is the Jacobian of the identity map?

\item What is the shape of the Jacobian for $F:\mathbb{R}^n \to \mathbb{R}^m$?

\item When is the Jacobian invertible?

\item What is the Jacobian determinant?

\item What is its geometric meaning?

\item How is it used in change of variables?

\item What does it mean when it is zero?

\item State the Inverse Function Theorem (informally).

\item What condition ensures local invertibility?

\end{enumerate}

\subsection*{Multiple Integrals}

\begin{enumerate}

\item What is a double integral of a function $f(x,y)$ over a region $R$?

\item How is the double integral defined using Riemann sums?

\item What is the geometric interpretation of a double integral?

\item What types of quantities can be computed using double integrals?

\item What is the difference between Type I and Type II regions?

\item How do we compute a double integral as an iterated integral?

\item State Fubini's theorem.

\item When can the order of integration be changed?

\item How is the area of a region computed using double integrals?

\item How is the average value of a function over a region defined?

\item What is a triple integral?

\item What is the geometric interpretation of a triple integral?

\item How can triple integrals be written as iterated integrals?

\item How is volume computed using triple integrals?

\item What are cylindrical coordinates?

\item What are spherical coordinates?

\item How do we convert between Cartesian and cylindrical coordinates?

\item How do we convert between Cartesian and spherical coordinates?

\item What is the Jacobian for cylindrical coordinates?

\item What is the Jacobian for spherical coordinates?

\item What is the change-of-variables formula for multiple integrals?

\item What is the geometric interpretation of the Jacobian determinant in integration?

\end{enumerate}
